% ==================================================================
%  徐森林、金亚东、薛春华《数学分析》（第三册）
%  复习题 15
%  源：原书印刷页 241-243（PDF 物理页 247-249）
%  底本：PaddleOCR-VL-1.6，见 chapters/_ch15_15_4_raw.md 末尾
%  约定：只用 exblock，不设 section（故不入目录）
%  全 13 题题面照录
% ==================================================================

\begin{exblock}{复习题 15}
\begin{enumerate}
  \item 证明: 对任何实数 $u$, 有
  \[
    \frac{1}{2\pi}\int_{0}^{2\pi}\mathrm{e}^{u\cos x}\cos(u\sin x)\,\mathrm{d}x=1.
\]

  \item 设 $D=\{(x,y)\mid 1\leqslant x^{2}+y^{2}\leqslant 4\}$, 计算积分
  \[
    \frac{1}{2\pi}\int_{0}^{+\infty}\left[\iint\limits_{D}\frac{\sin(t\sqrt{x^{2}+y^{2}})}{\sqrt{x^{2}+y^{2}}}\,\mathrm{d}x\mathrm{d}y\right]\mathrm{d}t.
\]

  \item 设函数 $f(x,u)$ 与 $\displaystyle\frac{\partial f(x,u)}{\partial u}$ 都在 $[a,+\infty)\times[\alpha,\beta]$ 上连续, 积分 $\displaystyle\int_{a}^{+\infty}\frac{\partial f(x,u)}{\partial u}\,\mathrm{d}x$ 在 $[\alpha,\beta]$ 上一致收敛, $\displaystyle\int_{a}^{+\infty}f(x,u)\,\mathrm{d}x$ 在某点 $u_{0}\in[\alpha,\beta]$ 收敛. 证明: 积分 $\displaystyle\int_{a}^{+\infty}f(x,u)\,\mathrm{d}x$ 在 $[\alpha,\beta]$ 上一致收敛于
  \[
    \varphi(u)=\int_{a}^{+\infty}f(x,u)\,\mathrm{d}x,
\]
  且
  \[
    \varphi'(u)=\frac{\mathrm{d}}{\mathrm{d}u}\int_{a}^{+\infty}f(x,u)\,\mathrm{d}x=\int_{a}^{+\infty}\frac{\partial f(x,u)}{\partial u}\,\mathrm{d}x.
\]

  \item 证明:
  \[
    \sum_{k=0}^{n}\mathrm{C}_{n}^{k}\frac{(-1)^{k}}{m+k+1}=\frac{n!m!}{(n+m+1)!}.
\]

  \item 设 $\displaystyle f(x)=\left(\int_{0}^{x}\mathrm{e}^{-t^{2}}\,\mathrm{d}t\right)^{2}$, $\displaystyle g(x)=\int_{0}^{1}\frac{\mathrm{e}^{-x^{2}(1+t^{2})}}{1+t^{2}}\,\mathrm{d}t$, 证明:
  \[
    f(x)+g(x)=\frac{\pi}{4}.
\]
  并由此推得 $\displaystyle\int_{0}^{+\infty}\mathrm{e}^{-x^{2}}\,\mathrm{d}x=\frac{\sqrt{\pi}}{2}$.

  \item 利用公式
  \[
    \Gamma(x)=\lim_{n\to+\infty}\frac{n^{x}n!}{x(x+1)\cdots(x+n)},
\]
  证明: 对任意实数 $x$ 有
  \[
    \sin x=x\prod_{n=1}^{\infty}\left(1-\frac{x^{2}}{n^{2}\pi^{2}}\right).
\]

  \item 按照下列步骤, 给出公式
  \[
    \mathrm{B}(p,q)=\frac{\Gamma(p)\Gamma(q)}{\Gamma(p+q)}
\]
  的另一个证明:

  (1)
  \[
    \Gamma(p)=2\int_{0}^{+\infty}u^{2p-1}\mathrm{e}^{-u^{2}}\,\mathrm{d}u.
\]

  (2)
  \[
    \Gamma(p)\Gamma(q)=\lim_{A\to+\infty}4\iint\limits_{[0,A]^{2}}u^{2p-1}v^{2q-1}\mathrm{e}^{-(u^{2}+v^{2})}\,\mathrm{d}u\mathrm{d}v.
\]

  (3)
  \[
    \begin{aligned}
      \lim_{A\to+\infty}\iint\limits_{D(A)}u^{2p-1}v^{2q-1}\mathrm{e}^{-(u^{2}+v^{2})}\,\mathrm{d}u\mathrm{d}v&=\frac{1}{4}\mathrm{B}(p,q)\Gamma(p+q),\\
      \lim_{A\to+\infty}\iint\limits_{D(\sqrt{2}A)}u^{2p-1}v^{2q-1}\mathrm{e}^{-(u^{2}+v^{2})}\,\mathrm{d}u\mathrm{d}v&=\frac{1}{4}\mathrm{B}(p,q)\Gamma(p+q),
    \end{aligned}
\]
  其中 $D(R)=\left\{(r,\theta)\mid 0\leqslant r\leqslant R, 0\leqslant\theta\leqslant\frac{\pi}{2}\right\}$.

  (4)
  \[
    \mathrm{B}(p,q)=\frac{\Gamma(p)\Gamma(q)}{\Gamma(p+q)}.
\]

  \item (1) 用两种不同方法证明:
  \[
    I(\alpha)=\int_{0}^{+\infty}\mathrm{e}^{-x^{2}}\cos(2\alpha x)\,\mathrm{d}x=\frac{\sqrt{\pi}}{2}\mathrm{e}^{-\alpha^{2}};
\]

  (2) 证明:
  \[
    J(\alpha)=\int_{0}^{+\infty}\mathrm{e}^{-x^{2}}\sin(2\alpha x)\,\mathrm{d}x=\mathrm{e}^{-\alpha^{2}}\int_{0}^{\alpha}\mathrm{e}^{-x^{2}}\,\mathrm{d}x.
\]

  \item (1) 证明:
  \[
    I(\alpha)=\int_{0}^{1}(1-x^{\alpha})^{-\frac{1}{\alpha}}\,\mathrm{d}x
\]
  当 $\alpha>1$ 时收敛; 当 $\alpha\leqslant 1$ 时发散.

  (2) 证明:
  \[
    I(\alpha)=\frac{\pi}{\alpha\sin\frac{\pi}{\alpha}},\quad\alpha>1.
\]

  \item 设 $\displaystyle f(x)=\int_{0}^{\frac{\pi}{2}}\ln(1-x^{2}\cos^{2}\theta)\,\mathrm{d}\theta$, $|x|<1$.

  (1) 用积分号下对 $x$ 求导, 然后积分得
  \[
    f(x)=\pi\ln\frac{1+\sqrt{1-x^{2}}}{2},\quad |x|<1;
\]

  (2) 对 $\displaystyle\ln(1-x^{2}\cos^{2}\theta)=-\sum_{n=1}^{\infty}\frac{x^{2n}}{n}\cos^{2n}\theta$ 以及幂级数的逐项积分证明:
  \[
    f(x)=-\frac{\pi}{2}\sum_{n=1}^{\infty}\frac{(2n-1)!!}{n\cdot(2n)!!}x^{2n};
\]

  (3) 根据 (1) 与 (2) 给出 $\ln(1+\sqrt{1-x^{2}})$ 的 Maclaurin 展开式, 以及 $f^{(k)}(0)$, $k=0,1,2,\cdots$.

  \item 用下列各方法证明:
  \[
    \int_{0}^{+\infty}\frac{\mathrm{e}^{-ax^{2}}-\mathrm{e}^{-bx^{2}}}{x^{2}}\,\mathrm{d}x=\sqrt{\pi}(\sqrt{b}-\sqrt{a}),\quad a>0,b>0.
\]

  (1) 应用积分号下求导;

  (2) 应用积分号下求积分;

  (3) 应用分部积分 $\displaystyle\int_{0}^{+\infty}(\mathrm{e}^{-ax^{2}}-\mathrm{e}^{-bx^{2}})\,\mathrm{d}\left(-\frac{1}{x}\right)$;

  (4) 化为二重积分
  \[
    \int_{0}^{+\infty}\frac{\mathrm{e}^{-ax^{2}}-\mathrm{e}^{-bx^{2}}}{x^{2}}\,\mathrm{d}x=\int_{0}^{+\infty}\frac{\mathrm{d}x}{x^{2}}\int_{ax^{2}}^{bx^{2}}\mathrm{e}^{-y}\,\mathrm{d}y=\iint\limits_{D}\frac{\mathrm{e}^{-y}}{x^{2}}\,\mathrm{d}x\mathrm{d}y,
\]
  其中 $D=\left\{(x,y)\mid 0\leqslant x<+\infty, ax^{2}\leqslant y\leqslant bx^{2}\right\}$;

  (5) 将
  \[
    \frac{\mathrm{e}^{-ax^{2}}-\mathrm{e}^{-bx^{2}}}{x^{2}}=\sum_{k=0}^{\infty}\frac{1}{(k+1)!}\mathrm{e}^{-bx^{2}}(b-a)^{k+1}x^{2k}
\]
  两边积分, 并由
  \[
    \int_{0}^{+\infty}\mathrm{e}^{-bx^{2}}\,\mathrm{d}x=\frac{1}{\sqrt{b}}\frac{\sqrt{\pi}}{2},\quad\int_{0}^{+\infty}\mathrm{e}^{-bx^{2}}x^{2k}\,\mathrm{d}x=\frac{\sqrt{\pi}}{2}\frac{1}{2}\frac{3}{2}\cdots\frac{2k-1}{2}\frac{1}{b^{\frac{2k+1}{2}}}
\]
  代入推得结论.

  \item 作变量代换 $\displaystyle t=\tan\frac{\varphi}{2}$ 与 $\displaystyle\sqrt{\frac{1-k}{1+k}}t=\tan\frac{\theta}{2}$, 证明:
  \[
    \begin{aligned}
      \int_{0}^{\pi}\left(\frac{\sin\varphi}{1+\cos\varphi}\right)^{\alpha-1}\frac{\mathrm{d}\varphi}{1+k\cos\varphi}&=\frac{1}{1+k}\left(\sqrt{\frac{1+k}{1-k}}\right)^{\alpha}\int_{0}^{\pi}\tan^{\alpha-1}\frac{\theta}{2}\,\mathrm{d}\theta\\
      &=\frac{1}{1+k}\left(\sqrt{\frac{1+k}{1-k}}\right)^{\alpha}\frac{\pi}{\sin\frac{\alpha}{2}\pi}.
    \end{aligned}
\]

  \item 设 $0<p,q<1$, 用下面两种方法证明:
  \[
    I(p,q)=\int_{0}^{+\infty}\frac{x^{p-1}-x^{q-1}}{(1+x)\ln x}\,\mathrm{d}x=\ln\left|\frac{\tan\frac{p\pi}{2}}{\tan\frac{q\pi}{2}}\right|.
\]

  (1) 利用 $\displaystyle\frac{\partial I}{\partial p}=\frac{\pi}{\sin p\pi}$ 及积分号下求导;

  (2) 由 $\displaystyle\left.\frac{x^{u-1}}{\ln x}\right|_{u=q}^{u=p}=\int_{q}^{p}x^{u-1}\,\mathrm{d}u$ 及积分号下求积分.
\end{enumerate}
\end{exblock}
